%%
%% This is file `satopdf.tex',
%% documentation for the satopdf package.
%%
%% Copyright (C) 2026 Shijin Chen
%%
%% This work may be distributed and/or modified under the conditions of the
%% LaTeX Project Public License, either version 1.3c of this license or (at
%% your option) any later version.
%%
%% Build: pdflatex satopdf && pdflatex satopdf
%%
\documentclass[a4paper]{ltxdoc}
\usepackage[T1]{fontenc}
\usepackage{lmodern}
\usepackage{booktabs}
\usepackage{array}
\usepackage{longtable}
\usepackage{enumitem}
\usepackage{xcolor}
\usepackage{hyperref}
\usepackage{satopdf}

\hypersetup{
  colorlinks=true,
  linkcolor=blue!50!black,
  urlcolor=blue!50!black,
  pdftitle={The satopdf package},
  pdfauthor={Shijin Chen},
  pdfsubject={Structural analysis diagrams for CAE to PDF (TikZ)},
}

\GetFileInfo{satopdf.sty}

\title{%
  The \textsf{satopdf} package%
  \thanks{This file describes version~\fileversion\ from \filedate.}%
}
\author{%
  Shijin Chen\\
  \texttt{sjchen1998@gmail.com}%
}
\date{Printed \today\\[0.5em] Package date: \filedate\ (v\fileversion)}

% Compact figure + optional caption text (standard doc style, not dual-column)
\newcommand{\STPfig}[1]{%
  \par\medskip
  \begin{center}#1\end{center}%
  \par\medskip}

\begin{document}
\maketitle

\begin{abstract}
The \textsf{satopdf} package is an independent TikZ-based library for drawing
\textbf{structural analysis diagrams} for materials mechanics and CAE
workflows (AutoCAD \texttt{LoadXData} $\rightarrow$ DXF import $\rightarrow$ PDF).
It provides a consistent 2D/3D command set under the prefix
\texttt{\textbackslash stp\ldots{}} for members (beam, truss, rod, fiber, ghost),
supports, hinges, loads, springs, internal-force envelopes, dimensioning,
welds, and shell faces.  CAE units are millimetres, newtons, and derived units
(N-mm, N$\cdot$mm, \ldots).
\end{abstract}

\tableofcontents
\clearpage

%==============================================================================
\section{Introduction}
%==============================================================================

\textsf{satopdf} draws publication-quality line figures of frames, trusses,
axial rods, and shell sketches that accompany CAE models.  Typical uses:

\begin{itemize}[nosep]
  \item Manual TikZ figures in lecture notes and reports;
  \item Auto-generated fragments from DXF via
        \texttt{dxf\_to\_satopdf2d.py} / \texttt{dxf\_to\_satopdf3d.py};
  \item Per-picture colour or black-and-white schemes.
\end{itemize}

\noindent\textbf{Licence.}
Copyright \textcopyright\ 2026 Shijin Chen.
Distributed under the LPPL~1.3c or later.
The Current Maintainer is Shijin Chen.
This work consists of \texttt{satopdf.sty} and \texttt{satopdf.pdf}
(documentation generated from \texttt{satopdf.tex}).

Control-sequence names never contain digits (TeX limitation).
User macros use the \texttt{stp}/\texttt{stpD} prefixes.

%==============================================================================
\section{Installation and requirements}
%==============================================================================

Place \texttt{satopdf.sty} on the \TeX{} search path
(e.g.\ \path{texmf/tex/latex/local/satopdf.sty}) and refresh the filename
database if required.  Dependencies:

\begin{itemize}[nosep]
  \item \LaTeXe{} with \textsf{xparse} (kernel);
  \item \textsf{tikz} with
        \texttt{calc}, \texttt{arrows.meta},
        \texttt{decorations.pathmorphing},
        \texttt{decorations.pathreplacing},
        \texttt{patterns}, \texttt{positioning},
        \texttt{backgrounds}, \texttt{3d};
  \item \textsf{amsmath}, \textsf{amssymb}, \textsf{ifthen}, \textsf{pgfkeys}.
\end{itemize}

\begin{verbatim}
\usepackage{satopdf}
\end{verbatim}

There are currently \textbf{no package options}.
Appearance is controlled by TikZ picture keys and \cs{stpset}.

%==============================================================================
\section{Getting started}
%==============================================================================

\subsection{Minimal 2D (millimetre canvas)}

\begin{verbatim}
\begin{tikzpicture}[x=1mm,y=1mm,stpScheme=color]
  \stpNode{A}{0,0}
  \stpNode{B}{1000,0}
  \stpNode{C}{2000,0}
  \stpBeam{A}{B}
  \stpBeam{B}{C}
  \stpSupport{A}[type=wall, angle=270]
  \stpSupport{C}[type=pin]
  \stpPointLoad{B}{0,-1}[$F$]
\end{tikzpicture}
\end{verbatim}

\STPfig{%
\begin{tikzpicture}[x=0.015mm,y=0.015mm,stpScheme=color]
  \stpNode{A}{0,0}
  \stpNode{B}{1000,0}
  \stpNode{C}{2000,0}
  \stpBeam{A}{B}
  \stpBeam{B}{C}
  \stpSupport{A}[type=wall, angle=270, size=8mm]
  \stpSupport{C}[type=pin, size=8mm]
  \stpset{force length=6mm}
  \stpPointLoad{B}{0,-1}[$F$]
\end{tikzpicture}}

\subsection{Scaled sketch coordinates}

\DescribeMacro{\stpScaling}
\cs{stpScaling}\marg{factor} multiplies subsequent \cs{stpPoint} inputs.
Default factor is~$1$ (TikZ centimetres).  Always set a scale when model
numbers are large, or coordinates are treated as centimetres.

\begin{verbatim}
\begin{tikzpicture}[stpScheme=color]
  \stpScaling{0.08}
  \stpPoint{A}{0}{0}
  \stpPoint{B}{40}{0}
  \stpBeam{A}{B}
  \stpSupport{A}[type=pin, size=8mm]
  \stpPointLoad{B}{0,-1}[$F$]
\end{tikzpicture}
\end{verbatim}

\subsection{Minimal 3D}

\begin{verbatim}
\stpDsetView{1}{-1}{1}
\renewcommand{\stpDscale}{0.08}
\begin{tikzpicture}[stpDCoords,stpScheme=color]
  \stpDpoint{A}{0}{0}{0}
  \stpDpoint{B}{28}{0}{0}
  \stpDmember{beam}{A}{B}
  \stpDaxis{A}
\end{tikzpicture}
\end{verbatim}

\STPfig{%
\stpDsetView{1}{-1}{1}%
\renewcommand{\stpDscale}{0.08}%
\begin{tikzpicture}[stpDCoords,stpScheme=color]
  \stpDpoint{A}{0}{0}{0}
  \stpDpoint{B}{28}{0}{0}
  \stpDmember{beam}{A}{B}
  \stpDaxis{A}
\end{tikzpicture}}

%==============================================================================
\section{Colour schemes and \texttt{\textbackslash stpset}}
\label{sec:scheme}
%==============================================================================

\subsection{Picture schemes}

On each \texttt{tikzpicture}, set \texttt{stpScheme=color} or
\texttt{stpScheme=bw} (aliases \texttt{stpColor}, \texttt{stpBW},
\texttt{mono}).  In black-and-white mode structural stroke colours become
black.  Soft-red \textbf{shell face fills} are unchanged.

\STPfig{%
\begin{tikzpicture}[stpScheme=color]
  \stpScaling{0.055}
  \stpPoint{A}{0}{0}\stpPoint{B}{40}{0}
  \stpBeam{A}{B}
  \stpSupport{A}[type=pin, size=8mm]
  \stpset{force length=6mm}
  \stpPointLoad{B}{0,-1}[$F$]
\end{tikzpicture}
\qquad
\begin{tikzpicture}[stpScheme=bw]
  \stpScaling{0.055}
  \stpPoint{A}{0}{0}\stpPoint{B}{40}{0}
  \stpBeam{A}{B}
  \stpSupport{A}[type=pin, size=8mm]
  \stpset{force length=6mm}
  \stpPointLoad{B}{0,-1}[$F$]
\end{tikzpicture}
\par\smallskip
{\footnotesize colour \qquad\qquad\qquad black-and-white}}

\subsection{\cs{stpset}}

\DescribeMacro{\stpset}
\cs{stpset}\marg{key=value,\ldots} adjusts defaults under \texttt{/stp}.

\begin{center}
\begin{tabular}{@{}llp{0.42\textwidth}@{}}
\toprule
Key & Typical & Effect \\
\midrule
\texttt{scaling} & \texttt{0.08} & Same as \cs{stpScaling} \\
\texttt{support size} & \texttt{12mm} & BC glyph base length \\
\texttt{hinge radius} & \texttt{1.5mm} & Hinge / rod-pin radius \\
\texttt{force length} & \texttt{10mm} & Point-load shaft \\
\texttt{lineload length} & \texttt{11mm} & Line-load arrow length \\
\texttt{lineload count} & \texttt{5} & Number of arrows \\
\texttt{moment radius} & \texttt{4mm} & Couple arc radius \\
\texttt{moment scale} & \texttt{0.015} & Envelope mm per $M$ unit \\
\texttt{moment segments} & \texttt{16} & Envelope polyline density \\
\texttt{deform scale} & \texttt{30} & Deformed-shape factor \\
\texttt{scheme} & \texttt{color}/\texttt{bw} & Prefer picture key \\
\texttt{dface shaft} & \texttt{0.55} & 3D face-load factor \\
\bottomrule
\end{tabular}
\end{center}

%==============================================================================
\section{Two-dimensional interface}
\label{sec:2d}
%==============================================================================

\subsection{Points}

\DescribeMacro{\stpPoint}
\cs{stpPoint}\marg{name}\marg{x}\marg{y} defines a coordinate multiplied by
the current \cs{stpScaling} factor.

\DescribeMacro{\stpNode}
\cs{stpNode}\marg{name}\marg{x,y} accepts CAE comma syntax.

\subsection{Members}

\DescribeMacro{\stpMember}
\cs{stpMember}\marg{type}\marg{A}\marg{B}\oarg{endA}\oarg{endB}

\begin{center}
\begin{tabular}{@{}llp{0.46\textwidth}@{}}
\toprule
Type & Stroke & End marks \\
\midrule
\texttt{beam}  & Thick black & \texttt{1} = solid joint dots \\
\texttt{fiber} & Thick + fibre & same as beam \\
\texttt{truss} & Medium & \texttt{0} = hollow hinges (default) \\
\texttt{rod}   & Medium dark blue & \texttt{0} = thin full pin circles;
  nonzero suppresses (use \cs{stpRodPin}) \\
\texttt{ghost} & Dashed & --- \\
\bottomrule
\end{tabular}
\end{center}

Shortcuts: \cs{stpBeam}, \cs{stpBeamFiber}, \cs{stpTruss}, \cs{stpRod},
\cs{stpGhost}.

\STPfig{%
\begin{tikzpicture}[stpScheme=color]
  \stpScaling{0.055}
  \stpPoint{A}{0}{0}\stpPoint{B}{35}{0}
  \stpPoint{C}{70}{0}\stpPoint{D}{105}{0}\stpPoint{E}{140}{0}
  \stpBeam{A}{B}
  \stpMember{fiber}{B}{C}
  \stpTruss{C}{D}
  \stpRod{D}{E}
  \stpNotation{node}{A}{beam}[below]
  \stpNotation{node}{B}{fiber}[below]
  \stpNotation{node}{C}{truss}[below]
  \stpNotation{node}{D}{rod}[below]
\end{tikzpicture}}

\paragraph{Rod end pins.}
\DescribeMacro{\stpRodPin}
\cs{stpRodPin}\marg{full\textbar gap}\marg{point}\oarg{beamA}\oarg{beamB}\oarg{rodOther}
\begin{itemize}[nosep]
  \item \texttt{full} --- thin hollow circle (default end mark);
  \item \texttt{gap} --- partial circle on the rod side of a multi-beam joint.
\end{itemize}

\STPfig{%
\begin{tikzpicture}[stpScheme=color]
  \stpScaling{0.055}
  \stpPoint{R1}{-80}{0}\stpPoint{J}{0}{0}\stpPoint{R2}{80}{0}
  \stpPoint{B1}{0}{50}\stpPoint{B2}{0}{-50}
  \stpMember{beam}{B1}{J}\stpMember{beam}{J}{B2}
  \stpMember{rod}{R1}{J}[0][1]\stpMember{rod}{J}{R2}[1][0]
  \stpRodPin{gap}{J}[B1][B2][R1]
  \stpSupport{R1}[type=pin, size=8mm]
  \stpSupport{R2}[type=roller, size=8mm]
\end{tikzpicture}}

\begin{verbatim}
\stpMember{rod}{R1}{J}[0][1]
\stpRodPin{gap}{J}[B1][B2][R1]
\end{verbatim}

\subsection{Supports and BC marks}

\DescribeMacro{\stpSupport}
\begin{verbatim}
\stpSupport{<node>}[type=..., angle=..., size=..., soft=..., fill=..., label=...]
\end{verbatim}
Types: \texttt{wall}, \texttt{fixed}, \texttt{pin}, \texttt{roller},
\texttt{rollerX}, \texttt{slider}, \texttt{spring}/\texttt{elastic},
\texttt{springCoil}, \texttt{torsion}, \texttt{torsionCW}.

\STPfig{%
\begin{tikzpicture}[stpScheme=color]
  \stpScaling{0.07}
  \stpPoint{A}{0}{0}\stpPoint{B}{25}{0}
  \stpPoint{C}{50}{0}\stpPoint{D}{75}{0}
  \stpPoint{E}{100}{0}\stpPoint{F}{125}{0}
  \stpGhost{A}{B}\stpGhost{B}{C}\stpGhost{C}{D}
  \stpGhost{D}{E}\stpGhost{E}{F}
  \stpSupport{A}[type=wall, size=8mm]
  \stpSupport{B}[type=fixed, size=8mm]
  \stpSupport{C}[type=pin, size=8mm]
  \stpSupport{D}[type=roller, size=8mm]
  \stpSupport{E}[type=slider, size=8mm]
  \stpSupport{F}[type=spring, size=8mm]
\end{tikzpicture}}

\DescribeMacro{\stpSupportFromMask}
\cs{stpSupportFromMask}\marg{UXUYUZRXRYRZ}\marg{node}\oarg{angle}
maps a six-bit CAE mask to pin, fixed, roller, or slider.

\STPfig{%
\begin{tikzpicture}[stpScheme=color]
  \stpScaling{0.065}
  \stpPoint{C}{0}{0}
  \stpPoint{T}{0}{30}
  \stpBeam{C}{T}
  \stpSupportFromMask{110000}{C}[0]
\end{tikzpicture}}

\DescribeMacro{\stpBCMark}
\cs{stpBCMark}\marg{node}\marg{dx,dy}\oarg{label} --- CAE two-circle glyph
along a unit direction (paper millimetres).

\STPfig{%
\begin{tikzpicture}[stpScheme=color]
  \stpPoint{N}{0}{0}
  \stpBCMark{N}{0,-1}[$BC$]
\end{tikzpicture}}

\subsection{Hinges}

\DescribeMacro{\stpHinge}
\cs{stpHinge}\marg{type}\marg{point}\oarg{\ldots}
Types: \texttt{full}, \texttt{gap}, \texttt{guide}, \texttt{axial},
\texttt{corner}.
For \texttt{gap}:
\cs{stpHinge}\texttt{\{gap\}\{J\}[beamA][beamB][trussOther]}
draws a half-hinge on the truss/rod side.

\STPfig{%
\begin{tikzpicture}[stpScheme=color]
  \stpScaling{0.065}
  \stpPoint{A}{0}{0}\stpPoint{B}{40}{0}
  \stpPoint{C}{70}{0}\stpPoint{T}{40}{35}
  \stpBeam{A}{B}\stpBeam{B}{C}\stpTruss{B}{T}
  \stpHinge{full}{A}
  \stpHinge{gap}{B}[A][C][T]
\end{tikzpicture}}

\subsection{Loads and springs}

\DescribeMacro{\stpLoad}
\cs{stpLoad}\marg{force\textbar moment\textbar momentCW}\marg{point}%
\oarg{rot}\oarg{len/sweep}\oarg{gap/radius}\oarg{label}
Radius arguments need an explicit unit (e.g.\ \texttt{13mm}); a bare
\texttt{4} is interpreted as $4\,\mathrm{cm}$.

\DescribeMacro{\stpPointLoad}
\cs{stpPointLoad}\marg{point}\marg{dx,dy}\oarg{label}

\DescribeMacro{\stpCouple}
\cs{stpCouple}\marg{point}\marg{cw\textbar ccw}\oarg{label}

\STPfig{%
\begin{tikzpicture}[stpScheme=color]
  \stpScaling{0.065}
  \stpPoint{A}{0}{0}\stpPoint{B}{45}{0}\stpPoint{C}{90}{0}
  \stpGhost{A}{B}\stpGhost{B}{C}
  \stpset{force length=6mm}
  \stpPointLoad{A}{0,-1}[$F$]
  \stpLoad{moment}{B}[0][200][13mm][$M$]
  \stpCouple{C}{ccw}[$M_c$]
\end{tikzpicture}}

\DescribeMacro{\stpLineLoad}
\cs{stpLineLoad}\marg{local\textbar global}\marg{A}\marg{B}%
\oarg{endA}\oarg{endB}\oarg{interval}\oarg{label}
draws a \textbf{double-rail} envelope (inner tip rail above the member,
outer head rail) with arrows between the rails.  Optional \texttt{endA}/
\texttt{endB} are centimetres added to the shaft; there are \emph{no}
tip circles on the member.

\STPfig{%
\begin{tikzpicture}[stpScheme=color]
  \stpScaling{0.065}
  \stpPoint{A}{0}{0}\stpPoint{B}{50}{0}
  \stpBeam{A}{B}
  \stpLineLoad{global}{A}{B}[0.35][0.35][0.2][$q$]
\end{tikzpicture}
\par\smallskip{\footnotesize\texttt{\textbackslash stpLineLoad\{global\}\ldots}}}

\DescribeMacro{\stpLineLoadCAE}
\cs{stpLineLoadCAE}\marg{A}\marg{B}\marg{dx,dy}\oarg{label}\oarg{count}
is the CAE vector form: tip circles sit on the member, a single rail
connects the arrow tails, and \texttt{count} sets the number of arrows
(default \cs{stpLineCount}).

\STPfig{%
\begin{tikzpicture}[stpScheme=color]
  \stpScaling{0.065}
  \stpPoint{A}{0}{0}\stpPoint{B}{50}{0}
  \stpBeam{A}{B}
  \stpLineLoadCAE{A}{B}{0,-1}[$q$][5]
\end{tikzpicture}
\par\smallskip{\footnotesize\texttt{\textbackslash stpLineLoadCAE}\ldots}}

\DescribeMacro{\stpSpringK}
\cs{stpSpringK}\marg{point}\marg{angle}\oarg{label}

\DescribeMacro{\stpSpringKR}
\cs{stpSpringKR}\marg{point}\oarg{label}
(arc uses a simple \texttt{->} tip).

\STPfig{%
\begin{tikzpicture}[stpScheme=color]
  \stpPoint{A}{0}{0}\stpPoint{B}{3}{0}
  \stpSpringK{A}{-90}[$K$]
  \stpSpringKR{B}[$K_R$]
\end{tikzpicture}}

\subsection{Internal forces and deformation}

\DescribeMacro{\stpMomentDiagram}
\cs{stpMomentDiagram}\oarg{qy=,scale=,label=,color=}\marg{A}\marg{B}\marg{Ms}\marg{Me}
builds
\[
  M(\xi)=M_s(1-\xi)+M_e\xi-\tfrac12 q_y L^2\xi(1-\xi),
\]
where $L$ is the drawn member length in millimetres and offsets are
$M\cdot\texttt{scale}$ in millimetres.  Verticals at the ends close the
envelope to the beam; they are not extra kinks in $M(x)$.
Diagram strokes use thin \texttt{stpMoment}/\texttt{stpShear} styles
(never beam weight).

\STPfig{%
\begin{tikzpicture}[stpScheme=color]
  \stpScaling{0.06}
  \stpPoint{A}{0}{0}\stpPoint{B}{40}{0}
  \stpBeam{A}{B}
  \stpMomentDiagram[qy=0.08,scale=1.2]{A}{B}{0}{0}
\end{tikzpicture}}

\DescribeMacro{\stpInternalForces}
\cs{stpInternalForces}\marg{A}\marg{B}\marg{vA}\marg{vB}\oarg{parabola}%
\oarg{color}\oarg{bendpos}
--- offsets in centimetres of paper.

\STPfig{%
\begin{tikzpicture}[stpScheme=color]
  \stpScaling{0.06}
  \stpPoint{A}{0}{0}\stpPoint{B}{40}{0}
  \stpBeam{A}{B}
  \stpInternalForces{A}{B}{0.55}{-0.4}[0][\stpClrShear]
\end{tikzpicture}}

\DescribeMacro{\stpDeform}
\cs{stpDeform}\marg{A}\marg{B}\marg{uxA,uyA}\marg{uxB,uyB}\oarg{scale}
displacements are multiplied by the deform scale (default~$30$).

\STPfig{%
\begin{tikzpicture}[stpScheme=color]
  \stpScaling{0.06}
  \stpPoint{A}{0}{0}\stpPoint{B}{40}{0}
  \stpBeam{A}{B}
  \stpDeform{A}{B}{0,0}{0,-0.12}[1]
\end{tikzpicture}}

\subsection{Dimensioning and notation}

\DescribeMacro{\stpDimensioning}
\cs{stpDimensioning}\marg{horiz\textbar vert\textbar parallel}\marg{A}\marg{B}%
\marg{offset}\oarg{label}
--- \texttt{horiz}/\texttt{vert} offsets are absolute TikZ coordinates.

\DescribeMacro{\stpDimX}\DescribeMacro{\stpDimY}
Envelope dimensions with optional paper $\Delta$ in parentheses.

\STPfig{%
\begin{tikzpicture}[stpScheme=color]
  \stpScaling{0.065}
  \stpPoint{SW}{0}{0}\stpPoint{SE}{40}{0}
  \stpPoint{NW}{0}{25}\stpPoint{NE}{40}{25}
  \stpGhost{SW}{SE}\stpGhost{SE}{NE}\stpGhost{NE}{NW}\stpGhost{NW}{SW}
  \stpDimX{SW}{SE}(0.35cm)[$L_x$]
  \stpDimY{SW}{NW}[$L_y$]
\end{tikzpicture}}

\DescribeMacro{\stpNotation}
\DescribeMacro{\stpTemperature}
\DescribeMacro{\stpInfluenceLine}
\DescribeMacro{\stpAddon}
\DescribeMacro{\stpAxisXY}
Auxiliary annotations.

\subsection{Welds and mass box}

\DescribeMacro{\stpSpotWeld}
\cs{stpSpotWeld}\marg{node}\oarg{radius}\oarg{label}

\DescribeMacro{\stpSeamWeld}
\cs{stpSeamWeld}\marg{A}\marg{B}\oarg{label}

\DescribeMacro{\stpSeamWeldPath}
\cs{stpSeamWeldPath}\marg{N1,N2,\ldots}\oarg{label}

Weld colour is \cs{stpClrWeld}.  Radii are paper millimetres.

\STPfig{%
\begin{tikzpicture}[stpScheme=color]
  \stpScaling{0.065}
  \stpPoint{A}{0}{0}\stpPoint{B}{40}{0}\stpPoint{C}{20}{25}
  \stpBeam{A}{B}\stpBeam{B}{C}
  \stpSpotWeld{B}[\stpWeldR][$w$]
  \stpSeamWeld{A}{B}
\end{tikzpicture}}

\DescribeMacro{\stpStructInfo}
\cs{stpStructInfo}\marg{$V=\ldots$}\marg{$m=\ldots$}
places a small box at the upper-right of the current picture bbox.

\DescribeMacro{\stpFace}
\cs{stpFace}\marg{N1,N2,\ldots} --- soft-red 2D plate fill
(scheme-independent).

%==============================================================================
\section{Three-dimensional interface}
\label{sec:3d}
%==============================================================================

\subsection{View and coordinates}

Call \textbf{before} \texttt{\textbackslash begin\{tikzpicture\}}:

\DescribeMacro{\stpDsetView}
\cs{stpDsetView}\marg{vx}\marg{vy}\marg{vz}\oarg{\ldots}
--- AutoCAD-style VPOINT.

\DescribeMacro{\stpDsetViewDir}
Alternative direction setter.

\DescribeMacro{\stpDsetCoords}
\cs{stpDsetCoords}\marg{x-ang}\marg{y-ang}\oarg{Dx}\oarg{Dy}\oarg{Dz}\oarg{z-ang}
--- legacy paper-angle oblique projection.

\DescribeMacro{\stpDsetAxis}
\cs{stpDsetAxis}\marg{xyz\textbar XYZ\textbar custom}\oarg{labels\ldots}

\DescribeMacro{\stpInitDimDelta}
Initialises paper $\Delta$ for envelope dimensions (run outside the
picture, before \texttt{tikzpicture}).

Picture options: \texttt{[stpDCoords,stpScheme=color]}.
Scale model units with \cs{stpDscale}.

\subsection{Points and members}

\DescribeMacro{\stpDpoint}
\cs{stpDpoint}\marg{name}\marg{x}\marg{y}\marg{z}

\DescribeMacro{\stpDnode}
\cs{stpDnode}\marg{name}\marg{x,y,z}

\DescribeMacro{\stpDmember}
\cs{stpDmember}\marg{beam\textbar truss\textbar rod\textbar ghost}\marg{A}\marg{B}%
\oarg{endA}\oarg{endB}

Shortcut: \cs{stpDrod}.

\DescribeMacro{\stpDrodPin}
\cs{stpDrodPin}\marg{full\textbar gap}\marg{point}\oarg{beamA}\oarg{beamB}\oarg{orient}

\DescribeMacro{\stpDaxis}
\cs{stpDaxis}\marg{point}\oarg{labX}\oarg{labY}\oarg{labZ}

\STPfig{%
\stpDsetView{1}{-1}{1}%
\renewcommand{\stpDscale}{0.08}%
\begin{tikzpicture}[stpDCoords,stpScheme=color]
  \stpDpoint{R1}{-10}{0}{0}\stpDpoint{J}{0}{0}{0}
  \stpDpoint{R2}{10}{0}{0}\stpDpoint{B1}{0}{12}{0}\stpDpoint{B2}{0}{-12}{0}
  \stpDmember{beam}{B1}{J}\stpDmember{beam}{J}{B2}
  \stpDmember{rod}{R1}{J}[0][1]\stpDmember{rod}{J}{R2}[1][0]
  \stpDrodPin{gap}{J}[B1][B2][0]
\end{tikzpicture}}

\subsection{Supports, hinges, loads, faces}

\DescribeMacro{\stpDsupport}
\cs{stpDsupport}\marg{pin\textbar wall\textbar spring\textbar fork\textbar fixed}%
\marg{point}\oarg{ux}\oarg{uy}\oarg{uz}\oarg{rx}\oarg{ry}\oarg{rz}

\DescribeMacro{\stpDhinge}
\cs{stpDhinge}\marg{full\textbar gap\textbar corner}\marg{point}\oarg{\ldots}
For \texttt{gap}, pass named nearby points (not raw \texttt{0,15,0} model
coordinates) so the clip path does not inflate the bounding box.

\STPfig{%
\stpDsetView{1}{-1}{1}%
\renewcommand{\stpDscale}{0.08}%
\begin{tikzpicture}[stpDCoords,stpScheme=color]
  \stpDpoint{A}{0}{0}{0}\stpDpoint{B}{28}{0}{0}
  \stpDpoint{G}{28}{6}{0}
  \stpDmember{beam}{A}{B}
  \stpDhinge{full}{A}
  \stpDhinge{gap}{B}[A][G][0]
\end{tikzpicture}}

\DescribeMacro{\stpDlineload}
\cs{stpDlineload}\marg{A}\marg{B}\oarg{dx,dy,dz}\oarg{count}\oarg{shaft}\oarg{label}
draws the same \textbf{tip rail + head rail} comb as \cs{stpLineLoad}
(not the CAE tip-circle style).

\DescribeMacro{\stpDpointload}
\DescribeMacro{\stpDface}
\cs{stpDface}\marg{N1,N2,\ldots} --- soft-red shell fill.

\DescribeMacro{\stpDfaceload}
Face pressure / traction glyph.

\STPfig{%
\stpDsetView{1}{-1}{1}%
\renewcommand{\stpDscale}{0.08}%
\begin{tikzpicture}[stpDCoords,stpScheme=color]
  \stpDpoint{A}{0}{0}{0}\stpDpoint{B}{30}{0}{0}
  \stpDmember{beam}{A}{B}
  \stpDlineload{A}{B}[0,0,-1][4][0.5][$q$]
\end{tikzpicture}
\par\smallskip{\footnotesize\texttt{\textbackslash stpDlineload} (tip+head rails)}}

\subsection{3D dimensions and welds}

\DescribeMacro{\stpDdimEnv}
\cs{stpDdimEnv}\marg{x\textbar y\textbar z}\marg{x0}\marg{y0}\marg{z0}\marg{L}%
\oarg{$\Delta$}\oarg{label}

\DescribeMacro{\stpDdimSeg}
\DescribeMacro{\stpDspotWeld}
\DescribeMacro{\stpDseamWeld}
\DescribeMacro{\stpDseamWeldPath}
Same roles as the 2D weld commands; follow the current view.

\DescribeMacro{\stpDstructInfo}
Alias of \cs{stpStructInfo}.

%==============================================================================
\section{Complete examples}
\label{sec:examples}
%==============================================================================

\subsection{2D SA template (\texttt{satopdf-frame.tex} / \texttt{2d\_example.dxf})}

The TikZ Studio 2D structural-analysis template (also written as
\texttt{satopdf-frame.tex}):

\STPfig{%
\begin{tikzpicture}[x=0.09mm,y=0.09mm,stpScheme=color]
  \stpset{support size=12mm, hinge radius=1.5mm}
  \stpNode{N0}{50.9619,0.0000}
  \stpNode{N1}{350.9619,0.0000}
  \stpNode{N2}{350.9619,116.6667}
  \stpNode{N3}{650.9619,0.0000}
  \stpNode{N4}{650.9619,233.3333}
  \stpNode{N5}{950.9619,0.0000}
  \stpNode{N6}{950.9619,350.0000}
  \stpNode{N7}{1250.9619,0.0000}
  \stpNode{N8}{1250.9619,180.0000}
  \stpMember{beam}{N6}{N5}
  \stpMember{beam}{N5}{N8}
  \stpMember{beam}{N5}{N7}
  \stpMember{beam}{N8}{N6}
  \stpMember{beam}{N4}{N3}
  \stpMember{beam}{N2}{N1}
  \stpMember{beam}{N0}{N1}
  \stpMember{beam}{N1}{N3}
  \stpMember{beam}{N3}{N5}
  \stpMember{beam}{N6}{N0}
  \stpMember{beam}{N5}{N4}
  \stpMember{beam}{N3}{N2}
  \stpMember{beam}{N8}{N7}
  \stpSupport{N0}[type=fixed, angle=0, size=12mm]
  \stpSupport{N7}[type=roller, angle=0, size=12mm]
  \stpLineLoad{global}{N6}{N0}[0.85][0.85][0.1667][$q=10$]
  \stpLoad{moment}{N8}[30][200][0][$M=1000$]
  \stpNode{DimSW}{50.9619,0.0000}
  \stpNode{DimSE}{1250.9619,0.0000}
  \stpNode{DimNW}{50.9619,350.0000}
  \stpDimX{DimSW}{DimSE}[$L_x=1200\,\mathrm{mm}$]
  \stpDimY{DimSW}{DimNW}[$L_y=350\,\mathrm{mm}$]
  \stpStructInfo{$V=8.72\times 10^{-3}\,\mathrm{m}^{3}$}{$m=68.5\,\mathrm{kg}$}
\end{tikzpicture}}

\begin{verbatim}
\begin{tikzpicture}[x=0.117mm,y=0.117mm,stpScheme=color]
  \stpset{support size=12mm, hinge radius=1.5mm}
  \stpNode{N0}{50.9619,0.0000}
  ...
  \stpMember{beam}{N6}{N0}
  \stpSupport{N0}[type=fixed, angle=0, size=12mm]
  \stpSupport{N7}[type=roller, angle=0, size=12mm]
  \stpLineLoad{global}{N6}{N0}[0.85][0.85][0.1667][$q=10$]
  \stpLoad{moment}{N8}[30][200][0][$M=1000$]
  \stpDimX{DimSW}{DimSE}[$L_x=1200\,\mathrm{mm}$]
  \stpStructInfo{$V=\ldots$}{$m=\ldots$}
\end{tikzpicture}
\end{verbatim}

\subsection{3D SA template (\texttt{satopdf-SA3D.tex} / \texttt{3d\_example.dxf})}

The TikZ Studio 3D structural-analysis template with
\texttt{\textbackslash stpDscale}=0.002:

\STPfig{%
% Drawing scale 1:5 (projected 480.8×334.8 cm @1:1 → A0 84.1×118.9 cm; ladder 1…5000)
\stpDsetView{1}{-1}{1}
\stpDsetAxis{xyz}
\renewcommand{\stpDscale}{0.002}
\renewcommand{\stpDsupLen}{0.35}
\renewcommand{\stpDhingeR}{0.9mm}
\renewcommand{\stpDaxLen}{1.2}
% Δ before tikzpicture (nullfont inside would zero char height)
\stpInitDimDelta
\begin{tikzpicture}[stpDCoords,stpScheme=color]
  \stpDpoint{N0}{0.0000}{0.0000}{0.0000}
  \stpDpoint{N1}{200.0000}{0.0000}{0.0000}
  \stpDpoint{N2}{600.0000}{0.0000}{0.0000}
  \stpDpoint{N3}{800.0000}{0.0000}{0.0000}
  \stpDpoint{N4}{0.0000}{1000.0000}{0.0000}
  \stpDpoint{N5}{0.0000}{1000.0000}{1200.0000}
  \stpDpoint{N6}{800.0000}{1000.0000}{0.0000}
  \stpDpoint{N7}{800.0000}{1000.0000}{1200.0000}
  \stpDpoint{N8}{0.0000}{2000.0000}{0.0000}
  \stpDpoint{N9}{0.0000}{2000.0000}{1200.0000}
  \stpDpoint{N10}{800.0000}{2000.0000}{0.0000}
  \stpDpoint{N11}{800.0000}{2000.0000}{1200.0000}
  \stpDpoint{N12}{0.0000}{3000.0000}{0.0000}
  \stpDpoint{N13}{0.0000}{3000.0000}{1200.0000}
  \stpDpoint{N14}{800.0000}{3000.0000}{0.0000}
  \stpDpoint{N15}{800.0000}{3000.0000}{1200.0000}
  \stpDpoint{N16}{0.0000}{4000.0000}{0.0000}
  \stpDpoint{N17}{0.0000}{4000.0000}{1200.0000}
  \stpDpoint{N18}{800.0000}{4000.0000}{0.0000}
  \stpDpoint{N19}{800.0000}{4000.0000}{1200.0000}
  \stpDpoint{N20}{0.0000}{5000.0000}{0.0000}
  \stpDpoint{N21}{0.0000}{5000.0000}{1200.0000}
  \stpDpoint{N22}{800.0000}{5000.0000}{0.0000}
  \stpDpoint{N23}{800.0000}{5000.0000}{1200.0000}
  \stpDpoint{N24}{0.0000}{6000.0000}{0.0000}
  \stpDpoint{N25}{200.0000}{6000.0000}{0.0000}
  \stpDpoint{N26}{600.0000}{6000.0000}{0.0000}
  \stpDpoint{N27}{800.0000}{6000.0000}{0.0000}
  \stpDaxis{N0}
  \stpDmember{beam}{N0}{N4}
  \stpDmember{beam}{N4}{N8}
  \stpDmember{beam}{N8}{N12}
  \stpDmember{beam}{N12}{N16}
  \stpDmember{beam}{N16}{N20}
  \stpDmember{beam}{N20}{N24}
  \stpDmember{beam}{N4}{N5}
  \stpDmember{beam}{N8}{N9}
  \stpDmember{beam}{N12}{N13}
  \stpDmember{beam}{N16}{N17}
  \stpDmember{beam}{N20}{N21}
  \stpDmember{truss}{N0}{N5}[1][1]
  \stpDmember{beam}{N9}{N13}
  \stpDmember{beam}{N13}{N17}
  \stpDmember{beam}{N17}{N21}
  \stpDmember{truss}{N21}{N24}[1][1]
  \stpDmember{beam}{N3}{N6}
  \stpDmember{beam}{N6}{N10}
  \stpDmember{beam}{N10}{N14}
  \stpDmember{beam}{N14}{N18}
  \stpDmember{beam}{N18}{N22}
  \stpDmember{beam}{N22}{N27}
  \stpDmember{beam}{N6}{N7}
  \stpDmember{beam}{N10}{N11}
  \stpDmember{beam}{N14}{N15}
  \stpDmember{beam}{N18}{N19}
  \stpDmember{beam}{N22}{N23}
  \stpDmember{truss}{N3}{N7}[1][1]
  \stpDmember{beam}{N11}{N15}
  \stpDmember{beam}{N15}{N19}
  \stpDmember{beam}{N19}{N23}
  \stpDmember{truss}{N23}{N27}[1][1]
  \stpDmember{beam}{N0}{N1}
  \stpDmember{beam}{N1}{N2}
  \stpDmember{beam}{N2}{N3}
  \stpDmember{beam}{N4}{N6}
  \stpDmember{beam}{N8}{N10}
  \stpDmember{beam}{N12}{N14}
  \stpDmember{beam}{N16}{N18}
  \stpDmember{beam}{N21}{N23}
  \stpDmember{beam}{N17}{N19}
  \stpDmember{beam}{N20}{N22}
  \stpDmember{beam}{N24}{N25}
  \stpDmember{beam}{N25}{N26}
  \stpDmember{beam}{N26}{N27}
  \stpDmember{beam}{N13}{N15}
  \stpDmember{truss}{N7}{N10}[1][1]
  \stpDmember{truss}{N10}{N15}[1][1]
  \stpDmember{truss}{N15}{N18}[1][1]
  \stpDmember{truss}{N18}{N23}[1][1]
  \stpDmember{truss}{N5}{N8}[1][1]
  \stpDmember{truss}{N8}{N13}[1][1]
  \stpDmember{truss}{N13}{N16}[1][1]
  \stpDmember{truss}{N16}{N21}[1][1]
  \stpDmember{beam}{N1}{N25}
  \stpDmember{beam}{N2}{N26}
  \stpDmember{beam}{N7}{N11}
  \stpDmember{beam}{N11}{N9}
  \stpDmember{beam}{N9}{N5}
  \stpDmember{beam}{N5}{N7}
  \stpDface{N21,N5,N7,N23}
  \stpDface{N23,N7,N6,N22}
  \stpDface{N21,N5,N4,N20}
  \stpDsupport{pin}{N24}[1][0][1][0][0][0]  % BC GLOBAL[101000]
  \stpDsupport{pin}{N27}[1][0][1][0][0][0]  % BC GLOBAL[101000]
  \stpDsupport{pin}{N0}[1][1][1][1][1][1]  % BC GLOBAL[111111]
  \stpDsupport{pin}{N3}[1][1][1][1][1][1]  % BC GLOBAL[111111]
  \stpDlineload{N2}{N26}[0,0,-1][6][0.55][$q=1.5$]
  \stpDlineload{N1}{N25}[0,0,-1][6][0.55][$q=1.5$]
  \stpDpointload{N5}{0,0,-1}[0.55][$F=1000\,\mathrm{N}$]
  \stpDfaceload{N21,N5,N7,N23}{0,0,-1}[0.55][$P_{\mathrm{face}}=-1\,\mathrm{MPa}$]
  \node[stpLabel, below left=1.5mm] at (N0) {$A$};
  \node[stpLabel, below right=1.5mm] at (N3) {$B$};
  \node[stpLabel, above left=1.5mm] at (N24) {$D$};
  \node[stpLabel, above right=1.5mm] at (N27) {$C$};
  \stpDdimEnv{x}{0.0000}{0.0000}{0.0000}{800.0000}[$L_x=800\,\mathrm{mm}$]
  \stpDdimEnv{y}{0.0000}{0.0000}{0.0000}{6000.0000}[$L_y=6000\,\mathrm{mm}$]
  \stpDdimEnv{z}{0.0000}{0.0000}{0.0000}{1200.0000}[$L_z=1200\,\mathrm{mm}$]
  % Structure mass properties (UR of current figure bbox)
  \stpStructInfo{$V=2.7\times 10^{-1}\,\mathrm{m}^{3}$}{$m=2.12\,\mathrm{t}$}
\end{tikzpicture}
}

\begin{verbatim}
\stpDsetView{1}{-1}{1}
\stpDsetAxis{xyz}
\renewcommand{\stpDscale}{0.002}
\stpInitDimDelta
\begin{tikzpicture}[stpDCoords,stpScheme=color]
  \stpDpoint{N0}{0}{0}{0}
  ...
  \stpDmember{beam}{N0}{N4}
  \stpDface{N21,N5,N7,N23}
  \stpDsupport{pin}{N0}[1][1][1][1][1][1]
  \stpDlineload{N1}{N25}[0,0,-1][6][0.55][$q=1.5$]
  \stpDfaceload{N21,N5,N7,N23}{0,0,-1}[0.55][$P_{\mathrm{face}}$]
  \stpDdimEnv{y}{0}{0}{0}{6000}[$L_y=6000\,\mathrm{mm}$]
  \stpStructInfo{$V=\ldots$}{$m=\ldots$}
\end{tikzpicture}
\end{verbatim}

\subsection{Distribution examples}

Further sources next to the package:

\begin{itemize}[nosep]
  \item \texttt{satopdf-frame.tex} --- 2D SA template (above);
  \item \texttt{satopdf-example-2d.tex} --- additional 2D gallery;
  \item \texttt{satopdf-SA3D.tex} --- 3D SA template (above);
  \item \texttt{satopdf-frame-3d.tex} --- additional 3D demo;
  \item \texttt{satopdf-face-demo.tex} --- shell faces.
\end{itemize}

%==============================================================================
\section{CAE / DXF conventions}
\label{sec:cae}
%==============================================================================

Python importers emit \textsf{satopdf} fragments from DXF entities tagged by
\texttt{LoadXData.lsp}:

\begin{center}
\begin{tabular}{@{}ll@{}}
\toprule
XData / concept & Drawing effect \\
\midrule
\texttt{MY\_CAE\_MEMBER\_KIND=ROD} &
  \texttt{\textbackslash stpMember\{rod\}} / \texttt{\textbackslash stpDmember\{rod\}} \\
\texttt{MY\_CAE\_SEC} & Section; rods drop orient/offset \\
\texttt{MY\_CAE\_LOAD} & Distributed $q$ (beams; span follows mesh) \\
\texttt{MY\_CAE\_POINT\_LOAD} & \cs{stpPointLoad} / \cs{stpDpointload} \\
\texttt{MY\_CAE\_BC} & Supports / \cs{stpDsupport} \\
\texttt{MY\_CAE\_SPOT\_WELD} / \texttt{SEAM} & Weld macros \\
Beam intersections \& mid-span welds &
  Split beams; \textbf{rods never split} \\
\bottomrule
\end{tabular}
\end{center}

Drawing heuristic (non-rod 3D only): axis-aligned segments draw as
\texttt{beam}; skewed as \texttt{truss}.  This is visual only, not an FE
classification.  Explicit \texttt{ROD} XData always selects \texttt{rod}.

%==============================================================================
\section{Units}
\label{sec:units}
%==============================================================================

\begin{itemize}[nosep]
  \item Geometry in the CAE pipeline: \textbf{mm};
  \item Forces: \textbf{N} (point); distributed often \textbf{N-mm}
        ($\equiv$~N/mm);
  \item Moments: \textbf{N$\cdot$mm};
  \item Springs: $K$ in N-mm, $K_R$ in Nmm/rad.
\end{itemize}

Paper sizes of hinges, weld marks, and tip circles are millimetres of
drawing paper and are independent of \cs{stpDscale}.

%==============================================================================
\section{Parameter defaults}
\label{sec:params}
%==============================================================================

\begin{center}
\begin{tabular}{@{}ll@{}}
\toprule
Macro / key & Default \\
\midrule
\cs{stpHugeLW} / member & $2.0\,\mathrm{pt}$ \\
\cs{stpTipR} & $1.35\,\mathrm{mm}$ \\
\cs{stpForceL} & $10\,\mathrm{mm}$ \\
\cs{stpHingeR} & $1.5\,\mathrm{mm}$ \\
\cs{stpMomScale} & $0.015$ \\
\cs{stpDefScale} & $30$ \\
\cs{stpDimOff} & $1.6\,\mathrm{cm}$ \\
\cs{stpDscale} & $1$ \\
\cs{stpFaceOpacity} & $0.45$ \\
\bottomrule
\end{tabular}
\end{center}

%==============================================================================
\section{Command summary}
\label{sec:summary}
%==============================================================================

{\small
\begin{longtable}{@{}p{0.42\textwidth}p{0.50\textwidth}@{}}
\toprule
Command & Role \\
\midrule
\endfirsthead
\toprule
Command & Role \\
\midrule
\endhead
\bottomrule
\endfoot
\cs{stpset} & Global keys \\
\cs{stpScaling} & 2D input scale \\
\cs{stpPoint}, \cs{stpNode} & 2D points \\
\cs{stpMember}, \cs{stpBeam}, \cs{stpTruss}, \cs{stpRod}, \ldots & 2D members \\
\cs{stpRodPin} & Rod end full/gap pins \\
\cs{stpSupport}, \cs{stpSupportFromMask}, \cs{stpBCMark} & 2D BC \\
\cs{stpHinge} & 2D hinges \\
\cs{stpLoad}, \cs{stpPointLoad}, \cs{stpCouple} & Point loads / couples \\
\cs{stpLineLoad}, \cs{stpLineLoadCAE} & Line loads \\
\cs{stpSpringK}, \cs{stpSpringKR} & Springs \\
\cs{stpInternalForces}, \cs{stpMomentDiagram}, \cs{stpDeform} & Diagrams \\
\cs{stpDimensioning}, \cs{stpDimX}, \cs{stpDimY} & Dimensions \\
\cs{stpSpotWeld}, \cs{stpSeamWeld}, \cs{stpSeamWeldPath} & 2D welds \\
\cs{stpFace}, \cs{stpStructInfo} & Faces / $V$/$m$ box \\
\cs{stpDsetView}, \cs{stpDsetCoords}, \cs{stpDsetAxis} & 3D view \\
\cs{stpInitDimDelta} & 3D dim $\Delta$ init \\
\cs{stpDpoint}, \cs{stpDnode}, \cs{stpDmember}, \cs{stpDrod} & 3D geometry \\
\cs{stpDrodPin} & 3D rod pins \\
\cs{stpDsupport}, \cs{stpDhinge} & 3D BC / hinges \\
\cs{stpDlineload}, \cs{stpDpointload}, \cs{stpDfaceload} & 3D loads \\
\cs{stpDface} & Shell faces \\
\cs{stpDdimEnv}, \cs{stpDdimSeg} & 3D dimensions \\
\cs{stpDspotWeld}, \cs{stpDseamWeld}, \ldots & 3D welds \\
\end{longtable}
}

%==============================================================================
\section{Change history}
\label{sec:changes}
%==============================================================================

\begin{description}
  \item[v0.4.13 (2026-11-21)]
        Rod end options \cs{stpRodPin}/\cs{stpDrodPin} (\texttt{full}/\texttt{gap});
        thin diagram strokes; BC-mark unit fix; documentation \texttt{satopdf.pdf}.
  \item[v0.4.12]
        Independent package header; author \& LPPL copyright;
        removed legacy third-party attribution comments.
  \item[v0.4.11]
        Welds, dim $\Delta$ hooks, rod member type foundations.
  \item[earlier]
        2D/3D core, supports, loads, shells, CAE contracts.
\end{description}

\end{document}
